% ECONOMICS_PROBLEM: {"id": "D63-59", "title": "EF1 and Pareto optimality for additive mixed manna", "jel_code": "D63", "source_id": "OP-0190"}
% !TeX program = xelatex
\documentclass[11pt,a4paper]{article}
\usepackage[margin=23mm]{geometry}
\usepackage{fontspec}
\setmainfont{Latin Modern Roman}
\usepackage{amsmath,amssymb,mathtools}
\usepackage{xcolor,enumitem,booktabs,longtable,array,xurl,hyperref}
\definecolor{muted}{HTML}{53636C}
\hypersetup{colorlinks=true,urlcolor=blue,linkcolor=blue}
\setlength{\parindent}{0pt}
\setlength{\parskip}{5pt}
\newcommand{\R}{\mathbb R}
\newcommand{\E}{\mathbb E}
\newcommand{\N}{\mathbb N}
\newcommand{\ind}{\mathbf 1}
\newcommand{\Var}{\operatorname{Var}}
\newcommand{\Cov}{\operatorname{Cov}}
\newcommand{\supp}{\operatorname{supp}}
\DeclareMathOperator*{\argmin}{arg\,min}
\DeclareMathOperator*{\argmax}{arg\,max}
\newcommand{\entrymeta}[3]{{\sffamily\small\color{muted}\textbf{\texttt{#1}}\quad|\quad #2\par #3\par}}
\newcommand{\Problemref}[1]{\texttt{#1}}
\newcommand{\JELref}[2]{\texttt{#2} (JEL \texttt{#1})}
\begin{document}
\section*{EF1 and Pareto optimality for additive mixed manna}
\textbf{Problem ID:} \texttt{D63-59}\quad \textbf{JEL:} \texttt{D63}\par
% BEGIN ECONOMICS STATEMENT
\label{problem:OP-0190}
% GAME_THEORY_PROBLEM: {"local_id": "GT-060", "original_number": 60, "kind": "P", "entry_kind": "statement_only", "published": false, "review_required": true, "category": "game_theory"}
\label{game:GT-060}
{\sffamily\small\color{muted}
\textbf{GT-060} | Fair division and efficiency | \textbf{P}: additive unweighted source-date question.
\par}
\subsection*{Definitions and mathematical statement}
Let $N=[n]$, $n\geq3$, $M$ finite, and $u_i(S)=\sum_{g\in S}u_{ig}$ for arbitrary real $u_{ig}$, permitting different signs for the same item across agents. A complete allocation $A$ partitions $M$. It is EF1 if, for every $i,j$, either $u_i(A_i)\geq u_i(A_j)$, or removing some item from $A_j$ yields $u_i(A_i)\geq u_i(A_j\setminus\{g\})$, or removing some item from $A_i$ yields $u_i(A_i\setminus\{g\})\geq u_i(A_j)$. It is Pareto optimal if there is no complete allocation $B$ with
\[
 \forall i,\quad u_i(B_i)\geq u_i(A_i),
 \qquad \exists k,\quad u_k(B_k)>u_k(A_k).
\]
Determine whether
\[
 \forall n\geq3\ \forall M\text{ finite}\ \forall(u_{ig})\in\mathbb R^{n\times|M|},
 \quad\exists A\text{ complete, EF1, and Pareto optimal}.
\]
\subsection*{Short English statement}
Can mixed goods and chores always be allocated both efficiently and envy-free up to one deletion?
\subsection*{Variants and scope boundaries}
This is an additive, unweighted question. Arbitrary set functions, weighted EF1, and fractional Pareto optimality are different targets; general-utility failures cannot be silently imported here. Polynomial computation would be a further question, even if existence is established.
\subsection*{Source relationship and status}
The explicit question is verified in the survey's Open Question 2. The catalogue's September 2026 reference could be confirmed through its indexed abstract but its alleged open-question passage could not be verified; it is not used to certify present-day openness. The missing additivity assumption in the original catalogue is restored.
\subsection*{References}
\begingroup\footnotesize\raggedright
\textbf{[S1]} Shengxin Liu, Xinhang Lu, Mashbat Suzuki, and Toby Walsh (2024 version). \emph{Mixed Fair Division: A Survey}. arXiv:2306.09564v4. \textit{Verified locator:} Definitions 3.1 and 3.3; Section 4.1, Open Question 2. \url{https://arxiv.org/html/2306.09564v4}
\par\textbf{[S2]} \emph{Weighted Fair Division of Indivisible Mixed Manna} (2026), arXiv:2609.01580. \textit{Verification limit:} indexed abstract establishes weighted/additive context, but the catalogue's asserted EF1+PO status passage is unverified. \url{https://arxiv.org/abs/2609.01580}

% Block audit notes for the editor (not an additional numbered section):
% 042: wrong original single-item citation; R two-item model.
% 044: original unrestricted DSIC ratio is 1 (VCG); source is protocol dominance.
%      October 2026 claimed GS resolution only indexed; full status unverified.
% 045: general 1/6 prophet bound solved in 2023; retain posted-item-price question.
% 048: source discusses receiver welfare; exact approximation constants editorial.
% 049: fixed two-good protocol editorial; one-good gains already ruled out.
% 050,052,055: complete editorial models for broad agendas, not verbatim conjectures.
% 051: cited paper resolves all alpha in [1/2,1]; H, no fabricated replacement.
% 053: algorithm vs IC separated; distance restrictions and expectation explicit.
% 054: universal truthfulness; planar scope and zero optimum convention fixed.
% 057: exact full-input algorithm strengthening labelled R; source abstract insufficient.
% 058: MMS-feasible is a valuation property, not instance-level MMS feasibility.
% 059: precisely fixes single-deletion either-side EF1 for arbitrary set functions.
% 060: restore additive assumption; recent source-open locator unverified.


% Formalized catalogue entries 61--77. Source versions checked 8 October 2026.
\par\endgroup

\subsection*{Common conventions: Source mathematical conventions}

\textbf{Finite games.} For a finite set $A$, $\Delta(A)$ is its probability
simplex. Players choose independent mixed strategies in Nash equilibrium;
correlation is allowed only when explicitly part of the solution concept.
In a game with payoff functions $u_i$ and strategy profile $x$, an additive
$\varepsilon$ Nash equilibrium satisfies
\[
 u_i(x)\geq u_i(x_i',x_{-i})-\varepsilon
 \quad\text{for every player $i$ and every admissible deviation $x_i'$}.
\]
For numerical approximation, payoffs are normalized as locally specified.
An $\varepsilon$ payoff residual is distinct from distance at most
$\varepsilon$ to the set of exact equilibria. Point convergence, convergence
of empirical play, and convergence in distance to an equilibrium set are
also distinct.

\textbf{Computational representations.} Unless stated otherwise, a complexity
question takes rational data with binary encoding; polynomial time is measured
in total bit length. A fixed-accuracy polynomial algorithm is not necessarily
polynomial in $1/\varepsilon$, and neither is necessarily polynomial in
$\log(1/\varepsilon)$. Succinct games and value, demand, or sampling oracles
must declare their interfaces. Exact real solutions require an explicit
representation; an unspecified list of arbitrary real numbers is not a
finite computational output. A claimed algorithmic barrier is meaningful
only for its specified model and reductions.

\textbf{Dynamic games.} Histories, information, observation, and allowable
randomization are part of the model. A stationary strategy depends only on
the current state; a positional strategy in a deterministic graph game
chooses one action at each controlled vertex. Nash equilibrium at an initial
state and equilibrium after every history are different requirements.
An expectation of a pathwise $\liminf$ payoff, a $\liminf$ of expected averages,
a limiting expected average, and a uniform finite-horizon equilibrium are
different criteria. None is silently substituted for another.

\textbf{Allocation and mechanisms.} A complete allocation assigns every item
exactly once unless otherwise stated. Zero-valued goods, negative values,
capacity constraints, feasibility, lotteries, and copies of goods affect
fairness definitions and are specified locally. Dominant-strategy incentive
compatibility, Bayesian incentive compatibility, universal truthfulness,
and truthfulness in expectation impose different inequalities. Whenever
an objective ratio has zero denominator, its convention is stated or those
instances are excluded.

\textbf{Infinite-dimensional models.} Expectations, integrals, stochastic
equations, and optimization objectives require their local regularity,
integrability, filtrations, and admissible domains. A backward--forward
mean-field system is a boundary-value problem; it is not an initial-value
semiflow without an additional construction. Long-time, large-population,
and vanishing-noise limits require an order of limits and a convergence
topology. If a minimum need not be attained, the objective is its infimum
or supremum and the approximation requirement is stated separately.
% END ECONOMICS STATEMENT
\end{document}
